%=====================================================================================
\section{Recognition audit}\label{sec:recognition}

A model that recognises a statement could, in principle, score it with knowledge of what happened next:
the market reaction, the commentary, the following decisions. The standard response is to hide the
statement's identity before scoring. This section tests whether that works for FOMC statements. It asks
two questions: can Fable identify a blinded or paraphrased statement when asked, and does it identify
statements without being asked while it scores them?

\subsection{Dating on request}

Fable is shown each blinded pre-2024 statement and told that months, years, names and events have been
replaced by placeholders. It is asked to identify the meeting, to give its best guess even if unsure, and
to end with the year, the month and a confidence level (Appendix~\ref{app:prompts}). The probe covers all
\R{probe.blind.n} pre-2024 files, including \R{corpus.prenotices} notices. The exact month identifies the
meeting except in the \R{corpus.multimonth} months that have two statements. The same probe is
run on a sample of \R{probe.para.n} paraphrased statements spread evenly over the period, with a note that
the text may have been reworded.

A model with no memory can still date statements roughly. Policy-rate levels, purchase amounts and the
boilerplate of each era change slowly, so neighbouring statements look alike. To measure how far the
text alone goes, the text-similarity baseline represents each blinded statement by the frequencies of its
words and word pairs (TF-IDF weights), finds the most similar \emph{other} blinded statement by cosine
similarity, and takes that statement's date. The baseline is generous: it sees the true date of every
other statement. It can never name the exact meeting, because its answer is always another meeting, so I
count the statements whose two nearest neighbours are the meetings just before and just after. For
those, a reader who knew the sequence of statements could in principle infer the exact meeting. This
count is an upper bound for this method, not for every reader without memory: someone who knows the
history of the policy rate, the purchase programmes and the guidance dates could date many statements
more precisely without having read them before. The baseline shows how far textual similarity alone
goes.

\begin{table}[t]
\centering
\caption{Dating blinded and paraphrased statements}\label{tab:dating}
\footnotesize
\setlength{\tabcolsep}{4pt}
% generated by paper/make_tables.py - do not edit
\begin{tabular}{@{}l cccc@{}}
\toprule
 & \multicolumn{3}{c}{Claude Fable 5.1, asked for the date} & Text-similarity \\
\cmidrule(lr){2-4}
 & Blinded & Paraphrased & Blinded & baseline \\
 & 2000--2023 & 2000--2023 (sample) & 2024--2026 & 2000--2023 \\
\midrule
Exact meeting & 202/206 (98.1\%) & 28/30 (93.3\%) & 19/22 (86.4\%) & $\le$ 89/206 (43.2\%)$^a$ \\
Within 1 month & 204/206 (99.0\%) & 29/30 (96.7\%) & 21/22 (95.5\%) & 106/206 (51.5\%) \\
Within 3 months & 206/206 (100.0\%) & 30/30 (100.0\%) & 22/22 (100.0\%) & 190/206 (92.2\%) \\
\addlinespace
ZLB statements, exact & 55/56 (98.2\%) & 6/8 (75.0\%) & -- & $\le$ 28/56 (50.0\%)$^a$ \\
Confidence ``high'' & 205/206 (99.5\%) & 29/30 (96.7\%) & 17/22 (77.3\%) & -- \\
\bottomrule
\end{tabular}

\tabnote{\emph{Notes:} Fable is asked for the year and month of the meeting. Paraphrased = an evenly
spaced sample of \R{probe.para.n} paraphrased statements. 2024--2026 = all blinded statements from 2024 on.
ZLB (zero-lower-bound) statements: January 2009 to October 2015, when the target range was the same at every
meeting (\R{corpus.zlbnotices} of them a notice). The 2000--2023 columns include \R{corpus.prenotices}
notices. Text-similarity baseline: each blinded statement takes the date of its most similar other
blinded statement (TF-IDF over words and word pairs, cosine). $^a$Statements whose two nearest neighbours
are the meetings just before and after, an upper bound on exact recovery for this method.}
\end{table}

\begin{figure}[t]
\centering
% generated by paper/make_tables.py
\begin{tikzpicture}
\begin{axis}[ybar, bar width=16pt, width=0.95\linewidth, height=6.2cm,
  ymin=0, ymax=112, ytick={0,25,50,75,100}, ylabel={Share of statements (\%)}, symbolic x coords={Exact meeting,Within 1 month,Within 3 months},
  xtick=data, enlarge x limits=0.2, legend style={at={(0.5,-0.14)}, anchor=north, legend columns=-1, font=\footnotesize, draw=none, /tikz/every even column/.append style={column sep=8pt}},
  legend cell align=left, nodes near coords, nodes near coords style={font=\scriptsize, /pgf/number format/fixed, /pgf/number format/precision=0}, ymajorgrids, major grid style={black!15}]
\addplot[fill=black!80, draw=black] coordinates {(Exact meeting,98.1) (Within 1 month,99.0) (Within 3 months,100.0)};
\addlegendentry{Fable, blinded (n=206)}
\addplot[fill=black!45, draw=black] coordinates {(Exact meeting,93.3) (Within 1 month,96.7) (Within 3 months,100.0)};
\addlegendentry{Fable, paraphrased (n=30)}
\addplot[pattern=north east lines, draw=black] coordinates {(Exact meeting,43.2) (Within 1 month,51.5) (Within 3 months,92.2)};
\addlegendentry{Text similarity, no memory (n=206)}
\end{axis}
\end{tikzpicture}

\caption{Dating accuracy. Fable on blinded and paraphrased statements against a text-similarity baseline
that has no memory. For the baseline, ``exact meeting'' is the upper bound described in the text.}\label{fig:dating}
\end{figure}

Table~\ref{tab:dating} and Figure~\ref{fig:dating} show the results. Fable names the exact month of
\R{probe.blind.exact} of the \R{probe.blind.n} blinded statements (\R{probe.blind.exactpct}) and places
every one within three months. The four misses (\R{probe.blind.misses}) are each an adjacent meeting. On
the zero-lower-bound statements, where the rate level gives no clue, it is exact on \R{probe.blind.zlb} of
\R{probe.blind.zlbn}. It reports high confidence on \R{probe.blind.high} statements. Its replies also
restore details that blinding removed: for the blinded statement of 18 December 2013 it names Eric
Rosengren as the dissenter, and for 25 June 2003 it says that Robert Parry dissented in favour of a larger
cut. Both are correct. Paraphrasing barely helps: Fable names the exact month of \R{probe.para.exact} of
the \R{probe.para.n} paraphrased statements.

The baseline shows what the text alone can do. It places \R{base.w3pct} of the statements within three
months and \R{base.w1pct} within one month, so dating to within a few months is not evidence of memory.
Exact identification is more telling. The baseline's upper bound on exact recovery is \R{base.brpct}
(\R{base.br} of \R{base.n}), and \R{base.zlb.br} of \R{probe.blind.zlbn} at the zero lower bound, against
Fable's \R{probe.blind.exactpct} and \R{probe.blind.zlb} of \R{probe.blind.zlbn}. The probe therefore shows
that Fable has knowledge sufficient to identify the individual meeting. Whether that knowledge is of the
statements themselves or of a detailed policy history cannot be separated here, but the
zero-lower-bound statements, which share a rate level and differ mainly in purchase amounts and wording,
are the hardest case for general knowledge, and the restored dissenters' names point to the statements
themselves.

\subsection{Recognition while scoring}

Dating on request shows that Fable \emph{can} identify the statements. A more direct question is whether
it does so when it is only asked for a score. I search every stored scoring reply for a date in the form
``$\langle$month$\rangle$ [day,] $\langle$year$\rangle$'' or ``$\langle$year$\rangle$
statement/meeting/decision'' and check whether it is the statement's own month and year. On blinded and
paraphrased text no date is in the input, so a correct date came from the model. On the original text of
the five-prompt runs, \R{nav.n} early statements contained their year in a navigation line
(Section~\ref{sec:data}), but not the month. The counts are lower bounds: most prompts ask for the number
only, and replies were stored cut at 300 characters.

\begin{table}[t]
\centering
\caption{Scoring replies that name the statement's date, unprompted}\label{tab:spontaneous}
\footnotesize
\setlength{\tabcolsep}{4pt}
% generated by paper/make_tables.py - do not edit
\begin{tabular}{@{}lll rrr@{}}
\toprule
Scores & Text & Prompt & Replies & Name a date & Own month and year \\
\midrule
Claude Fable 5.1, five prompts & original & all & 1228 & 245 (20.0\%) & 241 (19.6\%) \\
\addlinespace[2pt]
 &  & P1\_minimal & 255 & 21 (8.2\%) & 21 (8.2\%) \\
 &  & P2\_defined & 244 & 43 (17.6\%) & 43 (17.6\%) \\
 &  & P3\_analyst & 242 & 2 (0.8\%) & 2 (0.8\%) \\
 &  & P4\_reversed & 244 & 32 (13.1\%) & 32 (13.1\%) \\
 &  & P5\_reason\_first & 243 & 147 (60.5\%) & 143 (58.8\%) \\
\addlinespace
Claude Fable 5.1, P1 re-run & original & all & 206 & 20 (9.7\%) & 20 (9.7\%) \\
\addlinespace
Claude Fable 5.1, P1 & blinded & all & 206 & 23 (11.2\%) & 23 (11.2\%) \\
\addlinespace
Claude Fable 5.1, P1 & paraphrased & all & 206 & 7 (3.4\%) & 7 (3.4\%) \\
\addlinespace
Gemini 2.5 Flash, five prompts & original & all & 1140 & 0 (0.0\%) & 0 (0.0\%) \\
\addlinespace
GPT-4, P1 & original & all & 228 & 0 (0.0\%) & 0 (0.0\%) \\
\bottomrule
\end{tabular}

\tabnote{\emph{Notes:} A reply ``names a date'' if it contains ``$\langle$month$\rangle$ [day,]
$\langle$year$\rangle$'' or ``$\langle$year$\rangle$ statement/meeting/decision''. ``Own month and year'':
the named date is the statement's own. Replies stored before 3 October 2026 were cut at 300 characters,
so the shares are lower bounds. The five-prompt Fable file includes repeat runs of some cells.}
\end{table}

Table~\ref{tab:spontaneous} reports the counts. Across its five-prompt runs on the original text, Fable
names the statement's own month and year in \R{spont.scoresfable.all.pct} of its replies. The share
depends on whether the prompt lets it talk: \R{spont.scoresfable.p5.pct} for the reason-first prompt P5,
\R{spont.scoresfable.p3.pct} for the analyst prompt P3. On blinded text, with P1, it does so in
\R{spont.blindfable.all.pct} of replies, slightly more than the \R{spont.cleanfable.all.pct} of the same
prompt on the cleaned original. On paraphrased text the share falls to \R{spont.parafable.all.pct}.
Gemini never names a date in \R{spont.fullgemini.all.n} replies, nor does GPT-4 in
\R{spont.gpt4p1.all.n}.

\begin{table}[t]
\centering
\caption{Excerpts from Fable's scoring replies}\label{tab:quotes}
\footnotesize
% generated by paper/make_tables.py - do not edit
\begin{tabular}{@{}llp{10.2cm}@{}}
\toprule
Statement & Prompt & Excerpt of the scoring reply \\
\midrule
2013-10-30 (blinded) & P1\_minimal & \ldots{}This is the October 2013 FOMC statement: continued QE at \$85B/mo\ldots \\
2014-01-29 (blinded) & P1\_minimal & \ldots{}This is the January 2014 FOMC statement: a second measured taper\ldots \\
2000-12-19 (original) & P4\_reversed & \ldots{}inflation to economic weakness --- a clear dovish pivot signaling upcoming cuts (which came in January 2001)\ldots \\
2004-05-04 (original) & P2\_defined & \ldots{}removed at a pace that is likely to be measured'' --- a clear hawkish shift setting up the June 2004 hike,\ldots \\
2007-12-11 (original) & P5\_reason\_first & \ldots{}acknowledged intensifying housing correction and rising financial strains (markets had hoped for more)\ldots \\
2016-03-16 (original) & P1\_minimal & \ldots{}dissent (George). Overall mildly dovish relative to prior expectations (dot plot cut from 4 to 2 hikes)\ldots \\
2016-11-02 (original) & P5\_reason\_first & \ldots{}continued to strengthen'' --- explicit signal that a hike is approaching (widely read as teeing up December)\ldots \\
\bottomrule
\end{tabular}

\tabnote{\emph{Notes:} Verbatim excerpts from stored scoring replies (the blinded ones from P1 on blinded
text). Each was asked only for a stance score. Several cite information that was not in the statement:
later decisions, the market's reaction or commentary, or the projections released with it.}
\end{table}

Table~\ref{tab:quotes} gives examples. Some replies go beyond naming the meeting and bring in information
that was not in the statement: later decisions (cuts ``came in January 2001''; a statement was ``setting
up the June 2004 hike''), the market's reaction and later commentary (``markets had hoped for more'';
``widely read as teeing up December''), and the projections released alongside the statement (the
``dot plot cut from 4 to 2 hikes''). This is direct evidence that the model retrieves knowledge from
outside the text, including knowledge of later events, while it scores. It does not show that this knowledge moves
the score, but it shows that the scores cannot be assumed to come from the text alone.

\subsection{What recognition implies}

The same excerpts show a second habit: Fable often judges a statement relative to the previous one or
to what markets expected (``mildly dovish relative to prior expectations''). A score formed against
expectations is closer to a policy surprise than a reading of the text in isolation would be, and knowing
the expectations at each meeting is itself a form of recognition.

Recognition is necessary for a model to recall a statement's market reaction, but it is not evidence that
the model does so. The finding is therefore not that Fable's scores are contaminated. It is that the
tools meant to rule contamination out do not work on this corpus. Blinding and paraphrasing were meant to
create an unrecognised version of each statement, so that scores on recognised and unrecognised text
could be compared. With \R{probe.hold.recog} of the \R{probe.hold.n} probed hold meetings dated exactly,
there is no unrecognised group to compare with.

This finding complements \citet{glasserman2023}, who remove company names from news headlines to limit
look-ahead bias. Headlines about individual firms are many, short and mostly obscure. FOMC statements are
few, long, and widely read and discussed; their content identifies them. For documents like these,
anonymisation and paraphrase do not rule out look-ahead bias. Only models whose training data stop before the documents
\citep{he2025,drinkall2024}, or documents the model cannot have seen, can do that.

\subsection{Where Fable's knowledge ends}\label{sec:cutoff}

Running the same probe on the \R{probe.post.n} blinded statements from 2024 to 2026 locates the end of
Fable's knowledge (Appendix Table~\ref{tab:lateprobe}). It dates \R{probe.post.exact} of them exactly,
including every statement from the second half of 2024 and from 2025, all with high confidence. Its
confidence falls in 2026: of the four statements from the first half of 2026 it dates only one with high
confidence, and it misses July 2026. Its knowledge therefore extends to about early 2026. The 2024--2025
statements, often treated as a ``post-cutoff'' test sample, are not unseen for this model.
